\section{The model}
\label{sec:model}

\subsection{Scale and structure}

The model is a discrete simulation on a $128\times128$ grid of $16{,}384$
tiles, each carrying an integer elevation on 16 levels, a terrain type, a solar
exposure, a dust loading, and optionally a mineral deposit. One tick is one
simulated day. The implementation is approximately $12{,}800$ lines of
dependency-free JavaScript across 20 modules, governed by 52 named constants
collected in a single table so that every tunable quantity is visible in one
place.

The choice of a coarse integer grid over a continuous field is deliberate. The
model's purpose is to expose which constraints bind against each other, and
that structure survives discretisation; what discretisation buys is the ability
to run a complete settlement history in a few seconds, which is what makes the
sweeps in Section~\ref{sec:results} affordable.

\subsection{Illumination, and what follows from it}

Elevation is the primary field; everything else is derived. Each tile's solar
exposure is computed by casting rays outward in eight directions and testing
whether the local horizon occludes the Sun, with the horizon rise per unit
distance --- the parameter \texttt{sunSlope} --- set from the site's latitude.
Near the pole this parameter is small, so even shallow relief occludes, and
deep crater floors receive an exposure of exactly zero.

From that one field the model derives, in order:

\begin{itemize}
\item \textbf{Where ice is.} Ice is placed only where exposure falls below a
  shadow threshold, with richness increasing as exposure falls further. This
  reproduces the qualitative structure of polar cold traps
  \citep{paige2010diviner} without attempting their actual distribution.
\item \textbf{What power is available.} Photovoltaic output scales directly
  with local exposure, so the ground that holds ice generates nothing.
\item \textbf{What ground is worth.} Land value rises with exposure and with
  elevation, so shadowed floors are also the least valuable ground to build
  on.
\end{itemize}

This triple coupling is the model's central design decision. Because all three
follow from a single raycast, the antagonism between them is structural rather
than imposed: the most valuable resource sits on the least valuable ground, and
is hardest to power. No parameter tunes this relationship, and none can switch
it off.

\subsection{Networks, demand and progression}

Three utility networks must independently reach ground before it develops:
surface transit tubes, which serve by adjacency; power conduits, which
propagate a flood fill from generators and also conduct through developed
buildings; and buried atmosphere mains, which coexist with whatever occupies
the surface. Population and employment are driven by a three-index demand
model in which habitation grows while jobs outrun residents, and commerce and
industry grow while residents outrun jobs, with the two employment ratios
summing to unity so the system has a fixed point rather than a ceiling.

Colonies advance through four eras gated jointly on peak population and
accumulated research. The gate is conjunctive on purpose: population alone
would make research optional, and research alone would let a tiny well-funded
outpost build a skyline.
